\section{Bernard and the Militia of God}


\begin{quotex}
In 1926, Rene Guenon published an essay on the life of Saint Bernard of Clairvaux\footnote{\url{https://www.gornahoor.net/?page\_id=15638}}.
\end{quotex}


He obviously held Bernard in high regard as he demonstrates:

\begin{quotex}
Amongst the great figures of the Middle Ages, there are few whose study is more effective than that of Saint Bernard for the purpose of dissipating certain prejudices dear to the modern mind. What indeed could be more disconcerting for the modern mind than to see a pure contemplative --- one who always wished to be and to remain as such---called upon to play a dominant role in conducting the affairs of Church and State.
\end{quotex}


As a true contemplative, Bernard regarded profane philosophy as a lesser form of knowledge. That is why he opposed Abelard for relying solely on philosophical debate:

\begin{quotex}
Abelard, through his writings and teachings, had acquired for himself the reputation of being one of the most skillful dialecticians; he even made excessive use of dialectic, for instead of seeing in it only what it really is---a simple means for reaching an understanding of the truth---he regarded it almost as an end in itself, which tended to lead to an over-reliance on words ... Abelard did not know how to distinguish between what pertained to reason and what was higher than it; between profane philosophy and sacred wisdom.
\end{quotex}


In contrast, Guenon describes Bernard's mode of thinking like this:

\begin{quotex}
He resolved at a single blow the most arduous questions because his thinking did not proceed by means of a long series of discursive operations. What philosophers strove to reach by a circuitous route and by proceeding tentatively, he arrived at immediately, through \textbf{intellectual intuition}, a faculty without which no real metaphysics is possible, and without which one can only grasp a shadow of the truth.
\end{quotex}


Guenon accepts Renan's claim that critical thinking refutes the intellectual intuition. However, that just shows that critical thinking is itself refuted, since it cannot lead to any knowledge of the higher things. He elaborates:

\begin{quotex}
Thus the life of Saint Bernard could be seen as a refutation in advance of the errors of rationalism and pragmatism, errors considered to be opposed to each other, but in fact interdependent; at the same time, for those who examine his life impartially, it confounds and upturns all the preconceived ideas of ``scientific'' historians, who believe---with Renan --- that ``the negation of the supernatural constitutes the very essence of critical thinking'' --- a thesis with which we readily agree, but for the reason that we see in this incompatibility the exact opposite of what the moderns do, namely, a condemnation, not of the supernatural, but of ``critical thinking.''
\end{quotex}


Although miraculous healings were often associated with Bernard, he gave them only secondary importance.

\begin{quotex}
In the course of his journeys, Saint Bernard frequently reinforced his preaching by miraculous healings, which, for the crowds, were visible signs of his mission. Bernard himself was unwilling to speak of them. Perhaps he imposed this restriction on himself because of his great modesty; but he undoubtedly attributed only a secondary importance to these miracles, considering them simply a concession accorded by divine mercy to the weakness of the faith of the majority of the populace, in keeping with the words of Christ: ``Blessed are they that have not seen and yet have believed!''
\end{quotex}


This is what Guenon had to say about Bernard and the Templars:

\begin{quotex}
Bernard wrote the Rule for the Order of the Temple (Templars) and his commentary on the rule set forth, in terms of magnificent eloquence, the mission and the ideal of Christian chivalry, which he called the ``militia of God.'' These connections between the Abbot of Clairvaux and the \textbf{Order of the Temple}, which modern historians regard as merely a rather secondary episode in his life, assuredly had a completely different importance in the eyes of men of the Middle Ages; and we have shown elsewhere that these connections undoubtedly constitute the reason that Dante chose Saint Bernard as his guide in the highest circles of Paradise.
\end{quotex}


That demonstrates his connection with the esoteric current in the Middle Ages.

Bernard had a special devotion to the Virgin Mary.

\begin{quotex}
We must draw attention to a pre-eminent characteristic of Saint Bernard, namely, the central place which the cult of the Holy Virgin played in his life and in his writings. This produced a great flowering of legends, which may be the reason why Bernard has remained so popular. He loved to give the Holy Virgin the title of Our Lady, a usage which subsequently became generalized, doubtless due in large part to his influence; it is as if he were, as has been said, a true ``knight of Mary,'' and truly regarded her as his ``lady,'' in the chivalric sense of the word.
\end{quotex}


Guenon sums up Bernard's life like this:

\begin{quotex}
Having become a monk, Bernard always remained a knight, as did all those of his class; at the same time, one could say that he was in some way predestined to play the role of intermediary, conciliator, and arbiter between religious power and political power, since he combined in his person the nature of each. \emph{He was both monk and knight}: these two characteristics were those of the members of the ``militia of God,'' of the Order of the Temple; they were also, first and foremost, those of the author of their Rule, the great saint who was called the last of the Fathers of the Church, and whom some would regard as the prototype of Galahad, that perfect knight without blemish, the victorious hero of the quest for the \textbf{Holy Grail}.
\end{quotex}

See also \textit{Angelic States of Being}\footnote{See Section \ref{sec:AngelicStates} in this book.}.


\flrightit{Posted on 2022-01-25 by Cologero}

